An appendix is necessary if materials that relate to the thesis as a whole or to larger parts of it are difficult to include in the main body of the thesis. This is particularly the case if they would disrupt the overall coherence of the thesis due to their size. The content of the appendix must not contain anything that is necessary for understanding the text. It should not be used to expand the thesis and make it appear “thicker”.

The appendix is suitable for supplementary documents and materials, especially if these are difficult or impossible for the reader to access, e.g., unpublished company documents.

Particularly in the case of empirical work, the appendix can be used to identify the data sets, mathematical-statistical methods or programs used in more detail. 

If the appendix consists of several elements, the individual elements must be separated from each other by numbering.

If the appendix remains empty, this page can be deleted.
